% --- Archon physics formalization source begin ---
% archon:physics
% archon:covers PhyXMiniProblems/problem_phyx_mini_0038.lean
% archon:source-report reports/phyx_mini/problem_phyx_mini_0038.source.json

\chapter{Physics problem phyx\_mini\_0038}
\label{ch:PhyXMiniProblems_problem_phyx_mini_0038}

\paragraph{Problem source.}
\#\# Physical scenario

A concave mirror forms an image, on a wall 3.00 m in front of the mirror, of a headlamp filament 10.0 cm in front of the mirror.

\#\# Question

What is the focal length of the mirror?

\#\# Answer choices

- A: A: 8.68cm

- B: B: 9.68cm

- C: C: 9.66cm

- D: D: 9.05cm

\#\# Figure caption (auxiliary; use the image as primary evidence)

The image is a diagram illustrating ray tracing for a mirror with labeled components. Key elements include:

- **Mirror**: A concave mirror is on the right side.
- **Object**: Positioned on the right near the mirror; labeled with height \textbackslash{}( h = 5.00 \textbackslash{}, \textbackslash{}text\{mm\} \textbackslash{}).
- **Rays**: Drawn from the object, reflecting off the mirror, and converging at a point.
- **Screen**: On the left side, displaying the image.
- **Image**: Projected on the screen with unknown height \textbackslash{}( h' = ? \textbackslash{}).
- **Distances**: 
  - \textbackslash{}( s = 10.0 \textbackslash{}, \textbackslash{}text\{cm\} \textbackslash{}) (distance from object to mirror).
  - \textbackslash{}( s' = 3.00 \textbackslash{}, \textbackslash{}text\{m\} \textbackslash{}) (distance from mirror to screen).
- **Optic Axis**: A dashed line marking the principal axis of the system.
- **C**: Center of curvature of the mirror.
- **Radius**: \textbackslash{}( R = ? \textbackslash{}) (implies the radius of curvature is unknown).

The diagram represents the geometric relationships and parameters needed for understanding mirror optics.

\#\# Dataset metadata

Geometrical Optics; Physical Model Grounding Reasoning, Spatial Relation Reasoning

\paragraph{Recorded answer/context.}
B: B: 9.68cm

\paragraph{Figure/image path.}
/root/proposal\_for\_physic/science-mango/phyx\_mini\_run/phyx\_data/test\_image/38.png

\paragraph{Formalization target.}
create a compiling Lean file with sorry bodies at `PhyXMiniProblems/problem\_phyx\_mini\_0038.lean`. The Lean declarations must preserve the physical quantities, dimensions or dimensional roles, figure labels, governing-law hypotheses, and final relation expressed by this problem.
Use Mathlib/Physlib names found through LeanExplore where available. If a physics API is missing, introduce faithful local abstractions rather than scalar placeholder aliases.

\begin{theorem}[Physics formalization target]
\label{thm:physics:phyx_mini_0038:target}
\lean{PhyXMiniProblems.ProblemPhyXMini0038.problem_phyx_mini_0038}
\uses{def:physics:phyx-mini-0038:phyxminiproblems-problemphyxmini0038-centimetersvalue, def:physics:phyx-mini-0038:phyxminiproblems-problemphyxmini0038-concavemirrorimagingsetup, def:physics:phyx-mini-0038:phyxminiproblems-problemphyxmini0038-matchesfigurereadouts, def:physics:phyx-mini-0038:phyxminiproblems-problemphyxmini0038-hasphysicallengthsigns, def:physics:phyx-mini-0038:phyxminiproblems-problemphyxmini0038-obeyssphericalmirrorequation, def:physics:phyx-mini-0038:phyxminiproblems-problemphyxmini0038-obeysradiusfocallengthlaw, def:physics:phyx-mini-0038:phyxminiproblems-problemphyxmini0038-answerchoice, def:physics:phyx-mini-0038:phyxminiproblems-problemphyxmini0038-isnearesthundredthreadout}
The assigned autoformalize agent should translate this physics problem into Lean declarations in the covered file, with theorem and lemma proof bodies written as `by sorry`.
\end{theorem}
\begin{proof}
This is an autoformalization task, not a proof task. Produce statements that can later be proved by the physics prover without weakening the physical model.
\end{proof}

% --- Archon declaration topology begin ---
\section{Declaration topology}
\begin{definition}[Optical Length]
  \label{def:physics:phyx-mini-0038:phyxminiproblems-problemphyxmini0038-opticallength}
  \lean{PhyXMiniProblems.ProblemPhyXMini0038.OpticalLength}
  A signed physical length whose numerical value changes coherently with unit choice
\end{definition}
\begin{proof}
  This is the declared model definition of Optical Length.
\end{proof}
\begin{definition}[choices With Length Unit]
  \label{def:physics:phyx-mini-0038:phyxminiproblems-problemphyxmini0038-choiceswithlengthunit}
  \lean{PhyXMiniProblems.ProblemPhyXMini0038.choicesWithLengthUnit}
  Unit choices obtained from SI by changing only the length unit
\end{definition}
\begin{proof}
  This is the declared model definition of choices With Length Unit.
\end{proof}
\begin{definition}[length Readout]
  \label{def:physics:phyx-mini-0038:phyxminiproblems-problemphyxmini0038-lengthreadout}
  \lean{PhyXMiniProblems.ProblemPhyXMini0038.lengthReadout}
  \uses{def:physics:phyx-mini-0038:phyxminiproblems-problemphyxmini0038-opticallength, def:physics:phyx-mini-0038:phyxminiproblems-problemphyxmini0038-choiceswithlengthunit}
  The scalar readout of a physical length in the specified length unit
\end{definition}
\begin{proof}
  This is the declared model definition of length Readout.
\end{proof}
\begin{definition}[meters Value]
  \label{def:physics:phyx-mini-0038:phyxminiproblems-problemphyxmini0038-metersvalue}
  \lean{PhyXMiniProblems.ProblemPhyXMini0038.metersValue}
  \uses{def:physics:phyx-mini-0038:phyxminiproblems-problemphyxmini0038-opticallength}
  The scalar SI readout of a physical length, in metres
\end{definition}
\begin{proof}
  This is the declared model definition of meters Value.
\end{proof}
\begin{definition}[centimeters Value]
  \label{def:physics:phyx-mini-0038:phyxminiproblems-problemphyxmini0038-centimetersvalue}
  \lean{PhyXMiniProblems.ProblemPhyXMini0038.centimetersValue}
  \uses{def:physics:phyx-mini-0038:phyxminiproblems-problemphyxmini0038-opticallength, def:physics:phyx-mini-0038:phyxminiproblems-problemphyxmini0038-lengthreadout}
  The scalar readout of a physical length in centimetres
\end{definition}
\begin{proof}
  This is the declared model definition of centimeters Value.
\end{proof}
\begin{definition}[millimeters Value]
  \label{def:physics:phyx-mini-0038:phyxminiproblems-problemphyxmini0038-millimetersvalue}
  \lean{PhyXMiniProblems.ProblemPhyXMini0038.millimetersValue}
  \uses{def:physics:phyx-mini-0038:phyxminiproblems-problemphyxmini0038-opticallength, def:physics:phyx-mini-0038:phyxminiproblems-problemphyxmini0038-lengthreadout}
  The scalar readout of a physical length in millimetres
\end{definition}
\begin{proof}
  This is the declared model definition of millimeters Value.
\end{proof}
\begin{definition}[Spherical Mirror Kind]
  \label{def:physics:phyx-mini-0038:phyxminiproblems-problemphyxmini0038-sphericalmirrorkind}
  \lean{PhyXMiniProblems.ProblemPhyXMini0038.SphericalMirrorKind}
  Spherical-mirror kind as viewed from the incident-light side
\end{definition}
\begin{proof}
  This is the declared model definition of Spherical Mirror Kind.
\end{proof}
\begin{definition}[Mirror Image Nature]
  \label{def:physics:phyx-mini-0038:phyxminiproblems-problemphyxmini0038-mirrorimagenature}
  \lean{PhyXMiniProblems.ProblemPhyXMini0038.MirrorImageNature}
  Whether the reflected rays physically meet or only appear to meet
\end{definition}
\begin{proof}
  This is the declared model definition of Mirror Image Nature.
\end{proof}
\begin{definition}[Image Receiver]
  \label{def:physics:phyx-mini-0038:phyxminiproblems-problemphyxmini0038-imagereceiver}
  \lean{PhyXMiniProblems.ProblemPhyXMini0038.ImageReceiver}
  The surface on which the image is received in the physical scenario
\end{definition}
\begin{proof}
  This is the declared model definition of Image Receiver.
\end{proof}
\begin{definition}[Image Orientation]
  \label{def:physics:phyx-mini-0038:phyxminiproblems-problemphyxmini0038-imageorientation}
  \lean{PhyXMiniProblems.ProblemPhyXMini0038.ImageOrientation}
  Transverse orientation of the image relative to the object
\end{definition}
\begin{proof}
  This is the declared model definition of Image Orientation.
\end{proof}
\begin{definition}[Principal Axis Orientation]
  \label{def:physics:phyx-mini-0038:phyxminiproblems-problemphyxmini0038-principalaxisorientation}
  \lean{PhyXMiniProblems.ProblemPhyXMini0038.PrincipalAxisOrientation}
  Orientation of the principal optical axis displayed in the figure
\end{definition}
\begin{proof}
  This is the declared model definition of Principal Axis Orientation.
\end{proof}
\begin{definition}[Concave Mirror Imaging Setup]
  \label{def:physics:phyx-mini-0038:phyxminiproblems-problemphyxmini0038-concavemirrorimagingsetup}
  \lean{PhyXMiniProblems.ProblemPhyXMini0038.ConcaveMirrorImagingSetup}
  \uses{def:physics:phyx-mini-0038:phyxminiproblems-problemphyxmini0038-opticallength, def:physics:phyx-mini-0038:phyxminiproblems-problemphyxmini0038-sphericalmirrorkind, def:physics:phyx-mini-0038:phyxminiproblems-problemphyxmini0038-mirrorimagenature, def:physics:phyx-mini-0038:phyxminiproblems-problemphyxmini0038-imagereceiver, def:physics:phyx-mini-0038:phyxminiproblems-problemphyxmini0038-imageorientation, def:physics:phyx-mini-0038:phyxminiproblems-problemphyxmini0038-principalaxisorientation}
  The physical quantities and qualitative labels in the ray diagram. radiusOfCurvature is the axial distance from the mirror vertex to the point labelled C; objectHeight is the displayed h, and imageHeight is the unknown displayed h'. Object and image distances are positive distances in front of the mirror under the real-is-positive convention
\end{definition}
\begin{proof}
  This is the declared model definition of Concave Mirror Imaging Setup.
\end{proof}
\begin{definition}[Matches Figure Readouts]
  \label{def:physics:phyx-mini-0038:phyxminiproblems-problemphyxmini0038-matchesfigurereadouts}
  \lean{PhyXMiniProblems.ProblemPhyXMini0038.MatchesFigureReadouts}
  \uses{def:physics:phyx-mini-0038:phyxminiproblems-problemphyxmini0038-metersvalue, def:physics:phyx-mini-0038:phyxminiproblems-problemphyxmini0038-centimetersvalue, def:physics:phyx-mini-0038:phyxminiproblems-problemphyxmini0038-millimetersvalue, def:physics:phyx-mini-0038:phyxminiproblems-problemphyxmini0038-concavemirrorimagingsetup}
  Readouts and qualitative facts supplied by the problem and its figure: a concave mirror, a real inverted image on the wall/screen, s = 10.0 cm, s' = 3.00 m, and object height h = 5.00 mm on a horizontal optic axis. The unknown h', R, and focal length are deliberately not fixed here
\end{definition}
\begin{proof}
  This is the declared model definition of Matches Figure Readouts.
\end{proof}
\begin{definition}[Has Physical Length Signs]
  \label{def:physics:phyx-mini-0038:phyxminiproblems-problemphyxmini0038-hasphysicallengthsigns}
  \lean{PhyXMiniProblems.ProblemPhyXMini0038.HasPhysicalLengthSigns}
  \uses{def:physics:phyx-mini-0038:phyxminiproblems-problemphyxmini0038-metersvalue, def:physics:phyx-mini-0038:phyxminiproblems-problemphyxmini0038-concavemirrorimagingsetup}
  Positivity conditions for the length magnitudes in the depicted real-image configuration
\end{definition}
\begin{proof}
  This is the declared model definition of Has Physical Length Signs.
\end{proof}
\begin{definition}[Obeys Spherical Mirror Equation]
  \label{def:physics:phyx-mini-0038:phyxminiproblems-problemphyxmini0038-obeyssphericalmirrorequation}
  \lean{PhyXMiniProblems.ProblemPhyXMini0038.ObeysSphericalMirrorEquation}
  \uses{def:physics:phyx-mini-0038:phyxminiproblems-problemphyxmini0038-concavemirrorimagingsetup}
  The paraxial spherical-mirror equation 1/f = 1/s + 1/s', expressed in the division-free, dimensionally homogeneous form f * (s + s') = s * s' in every unit system
\end{definition}
\begin{proof}
  This is the declared model definition of Obeys Spherical Mirror Equation.
\end{proof}
\begin{definition}[Obeys Radius Focal Length Law]
  \label{def:physics:phyx-mini-0038:phyxminiproblems-problemphyxmini0038-obeysradiusfocallengthlaw}
  \lean{PhyXMiniProblems.ProblemPhyXMini0038.ObeysRadiusFocalLengthLaw}
  \uses{def:physics:phyx-mini-0038:phyxminiproblems-problemphyxmini0038-concavemirrorimagingsetup}
  For a paraxial spherical mirror, the radius of curvature is twice its focal length
\end{definition}
\begin{proof}
  This is the declared model definition of Obeys Radius Focal Length Law.
\end{proof}
\begin{definition}[Answer Choice]
  \label{def:physics:phyx-mini-0038:phyxminiproblems-problemphyxmini0038-answerchoice}
  \lean{PhyXMiniProblems.ProblemPhyXMini0038.AnswerChoice}
  Labels of the four displayed multiple-choice answers
\end{definition}
\begin{proof}
  This is the declared model definition of Answer Choice.
\end{proof}
\begin{definition}[focal Length Centimeters]
  \label{def:physics:phyx-mini-0038:phyxminiproblems-problemphyxmini0038-answerchoice-focallengthcentimeters}
  \lean{PhyXMiniProblems.ProblemPhyXMini0038.AnswerChoice.focalLengthCentimeters}
  \uses{def:physics:phyx-mini-0038:phyxminiproblems-problemphyxmini0038-answerchoice}
  The focal-length readout printed beside each answer choice, in centimetres
\end{definition}
\begin{proof}
  This is the declared model definition of focal Length Centimeters.
\end{proof}
\begin{definition}[Is Nearest Hundredth Readout]
  \label{def:physics:phyx-mini-0038:phyxminiproblems-problemphyxmini0038-isnearesthundredthreadout}
  \lean{PhyXMiniProblems.ProblemPhyXMini0038.IsNearestHundredthReadout}
  reported is a nearest-hundredth decimal readout of exact: it is an integer number of hundredths and differs from the exact value by at most half a hundredth
\end{definition}
\begin{proof}
  This is the declared model definition of Is Nearest Hundredth Readout.
\end{proof}
\begin{lemma}[focal Length Centimeters eq three Hundred div thirty One]
  \label{lem:physics:phyx-mini-0038:phyxminiproblems-problemphyxmini0038-focallengthcentimeters-eq-threehundred-div-thirtyone}
  \lean{PhyXMiniProblems.ProblemPhyXMini0038.focalLengthCentimeters_eq_threeHundred_div_thirtyOne}
  \uses{def:physics:phyx-mini-0038:phyxminiproblems-problemphyxmini0038-centimetersvalue, def:physics:phyx-mini-0038:phyxminiproblems-problemphyxmini0038-concavemirrorimagingsetup, def:physics:phyx-mini-0038:phyxminiproblems-problemphyxmini0038-matchesfigurereadouts, def:physics:phyx-mini-0038:phyxminiproblems-problemphyxmini0038-obeyssphericalmirrorequation}
  The mirror equation and the two stated axial distances determine the exact focal length as 300/31 cm
\end{lemma}
\begin{proof}
  The conclusion follows from the governing relations and figure readouts named in the statement of focal Length Centimeters eq three Hundred div thirty One.
\end{proof}
% --- Archon declaration topology end ---
% --- Archon physics formalization source end ---
